% =====================================================================
%  Thème 2 — Chasles et combinaisons — Niveau FACILE — Corrigé
% =====================================================================
\documentclass[11pt,a4paper]{article}
\usepackage[utf8]{inputenc}
\usepackage[T1]{fontenc}
\usepackage{lmodern}
\usepackage[french]{babel}
\usepackage{amsmath,amssymb,mathtools}
\usepackage{xcolor}
\usepackage{enumitem}
\usepackage{fancyhdr}
\usepackage[a4paper,margin=2.1cm,headheight=26pt,headsep=10pt]{geometry}

\definecolor{primary}{RGB}{222,70,80}
\definecolor{primarydark}{RGB}{150,30,42}
\definecolor{excol}{RGB}{0,135,90}

\pagestyle{fancy}\fancyhf{}
\fancyhead[L]{\small\textcolor{primarydark}{\textbf{Fiche 2} — Calcul vectoriel}}
\fancyhead[R]{\small\textcolor{primarydark}{\textsc{Thème 2 — Facile — Corrigé}}}
\fancyfoot[C]{\small\thepage}
\renewcommand{\headrulewidth}{0.8pt}

\newcommand{\vc}[1]{\overrightarrow{#1}}
\providecommand{\norm}[1]{\left\lVert #1\right\rVert}
\newcommand{\cor}[1]{\medskip\noindent{\color{primarydark}\bfseries\large Corrigé #1.}\par\smallskip}
\newcommand{\rep}{\textcolor{excol}{$\blacktriangleright$}\ }

\begin{document}
\begin{center}
{\LARGE\bfseries\color{primarydark}Chasles et combinaisons — Corrigé Facile}\\[2pt]
{\color{primary}\rule{0.5\linewidth}{1.1pt}}
\end{center}
\vspace{2mm}

\cor{1}
\begin{enumerate}[label=\textbf{\alph*)},leftmargin=2em,itemsep=2pt]
  \item $\vc{AB}+\vc{BC}=\vc{AC}$.
  \item $\vc{MN}+\vc{NP}=\vc{MP}$.
  \item $\vc{RS}+\vc{ST}+\vc{TU}=\vc{RU}$.
  \item $\vc{EF}+\vc{FG}+\vc{GH}+\vc{HK}=\vc{EK}$.
\end{enumerate}
\rep La lettre intermédiaire commune s'efface à chaque étape ; il reste la première et la dernière.

\cor{2}
\begin{enumerate}[label=\textbf{\alph*)},leftmargin=2em,itemsep=2pt]
  \item $\vc{AB}+\vc{BA}=\vc{AA}=\vec0$.
  \item $\vc{AB}+\vc{BC}+\vc{CA}=\vc{AC}+\vc{CA}=\vc{AA}=\vec0$.
  \item $-\vc{DE}=\vc{ED}$.
  \item $\vc{PQ}+\vc{QP}+\vc{PQ}=\vec0+\vc{PQ}=\vc{PQ}$.
\end{enumerate}
\rep Un circuit fermé (on revient au point de départ) donne toujours le vecteur nul.

\cor{3}
\begin{enumerate}[label=\textbf{\alph*)},leftmargin=2em,itemsep=2pt]
  \item $\vc{AB}-\vc{CB}=\vc{AB}+\vc{BC}=\vc{AC}$.
  \item $\vc{AC}-\vc{AB}=\vc{AC}+\vc{CA}\;$? Non : on écrit $-\vc{AB}=\vc{BA}$, donc
        $\vc{AC}-\vc{AB}=\vc{BA}+\vc{AC}=\vc{BC}$.
  \item $\vc{MP}-\vc{MN}=\vc{NM}+\vc{MP}=\vc{NP}$.
\end{enumerate}
\rep Règle utile : $\vc{XA}-\vc{XB}=\vc{BA}$ (les origines communes $X$ « disparaissent »,
il reste \og{}extrémité de gauche moins celle de droite\fg{}). Ici b) : $\vc{AC}-\vc{AB}=\vc{BC}$.

\cor{4}
\begin{enumerate}[label=\textbf{\alph*)},leftmargin=2em,itemsep=2pt]
  \item $\vc{AC}=\vc{AB}+\vc{BC}$.
  \item $\vc{AD}=\vc{AB}+\vc{BC}+\vc{CD}$.
  \item $\vc{MN}=\vc{MO}+\vc{ON}$ (valable pour \emph{n'importe quel} point $O$).
\end{enumerate}
\rep On peut toujours \emph{insérer} un point : c'est Chasles lu de droite à gauche.

\cor{5}
\begin{enumerate}[label=\textbf{\alph*)},leftmargin=2em,itemsep=2pt]
  \item $\vc{AB}=\vc{DC}$ (côtés opposés du parallélogramme).
  \item $\vc{AD}=\vc{BC}$.
  \item $\vc{AB}+\vc{AD}=\vc{AC}$ (diagonale issue de $A$).
\end{enumerate}
\rep Dans $ABCD$ pris dans l'ordre, les côtés opposés sont des vecteurs \emph{égaux} :
$\vc{AB}=\vc{DC}$ et $\vc{AD}=\vc{BC}$.

\end{document}
